% GENERATED FILE. Edit canonical lessons, then run npm run generate:books.
% canonical-source-hash: fnv1a64:d1c50e89e391f379
% canonical-lessons: SA-C143-pratyagacchati, SA-C143-sete, SA-C143-aste, SA-C143-bibheti, SA-C143-dhyayati

\chapter{Comes back, Lies down, Sits, Is afraid, Meditates}
\label{ch:sa-comes-back-lies-down-sits-is-afraid-meditates}
\hlchaptermodality{143}

\begin{chapteropening}
\textbf{By the end of this chapter:} \emph{I can say comes back, lies down, sits, is afraid and meditates.}

Complete the last lesson of chapter 143: I can say comes back, lies down, sits, is afraid and meditates.
\end{chapteropening}

\section[\texorpdfstring{\sk{प्रत्यागच्छति} (pratyāgacchati)}{प्रत्यागच्छति (pratyāgacchati)}]{\sk{प्रत्यागच्छति} (pratyāgacchati) — comes back, returns}

\label{lesson:SA-C143-pratyagacchati}



\begin{quote}
Before the new one: say the Sanskrit for shows, then the Sanskrit for worships.
\end{quote}

\subsection*{You'll want to know: \sk{प्रत्यागच्छति}}
\textbf{\sk{प्रत्यागच्छति}} — \emph{pratyāgacchati} — \textquotedblleft{}comes back, returns\textquotedblright{}.

\begin{grammarlens}[title={What you've built}]
One more word for this chapter.
\end{grammarlens}

\subsection*{Your turn}
\begin{itemize}
\raggedright
  \item \emph{Say it:} \emph{pratyāgacchati}
  \item \emph{Say it:} \emph{pratyāgacchati}, once more
  \item \emph{Say it:} say \emph{pratyāgacchati}
  \item \emph{Recall:} say the Sanskrit for shows, then the Sanskrit for worships, then say \emph{pratyāgacchati} again
\end{itemize}

\begin{tcolorbox}[breakable,colback=teal!4,colframe=teal!35!black,title={Before you move on}]
What does \sk{प्रत्यागच्छति} mean? (\textquotedblleft{}Comes back, returns\textquotedblright{}.) Say it once more.
\end{tcolorbox}
\section[\texorpdfstring{\sk{शेते} (śete)}{शेते (śete)}]{\sk{शेते} (śete) — lies down}

\label{lesson:SA-C143-sete}



\begin{quote}
Before the new one: say the Sanskrit for worships, then the Sanskrit for comes back.
\end{quote}

\subsection*{You'll want to know: \sk{शेते}}
\textbf{\sk{शेते}} — \emph{śete} — \textquotedblleft{}lies down\textquotedblright{}.

\begin{grammarlens}[title={What you've built}]
One more word for this chapter.
\end{grammarlens}

\subsection*{Your turn}
\begin{itemize}
\raggedright
  \item \emph{Say it:} \emph{śete}
  \item \emph{Say it:} \emph{śete}, once more
  \item \emph{Say it:} say \emph{śete}
  \item \emph{Recall:} say the Sanskrit for worships, then the Sanskrit for comes back, then say \emph{śete} again
\end{itemize}

\begin{tcolorbox}[breakable,colback=teal!4,colframe=teal!35!black,title={Before you move on}]
What does \sk{शेते} mean? (\textquotedblleft{}Lies down\textquotedblright{}.) Say it once more.
\end{tcolorbox}
\section[\texorpdfstring{\sk{आस्ते} (āste)}{आस्ते (āste)}]{\sk{आस्ते} (āste) — sits, stays}

\label{lesson:SA-C143-aste}



\begin{quote}
Before the new one: say the Sanskrit for comes back, then the Sanskrit for lies down.
\end{quote}

\subsection*{You'll want to know: \sk{आस्ते}}
\textbf{\sk{आस्ते}} — \emph{āste} — \textquotedblleft{}sits, stays\textquotedblright{}.

\begin{grammarlens}[title={What you've built}]
One more word for this chapter.
\end{grammarlens}

\subsection*{Your turn}
\begin{itemize}
\raggedright
  \item \emph{Say it:} \emph{āste}
  \item \emph{Say it:} \emph{āste}, once more
  \item \emph{Say it:} say \emph{āste}
  \item \emph{Recall:} say the Sanskrit for comes back, then the Sanskrit for lies down, then say \emph{āste} again
\end{itemize}

\begin{tcolorbox}[breakable,colback=teal!4,colframe=teal!35!black,title={Before you move on}]
What does \sk{आस्ते} mean? (\textquotedblleft{}Sits, stays\textquotedblright{}.) Say it once more.
\end{tcolorbox}
\section[\texorpdfstring{\sk{बिभेति} (bibheti)}{बिभेति (bibheti)}]{\sk{बिभेति} (bibheti) — is afraid}

\label{lesson:SA-C143-bibheti}



\begin{quote}
Before the new one: say the Sanskrit for lies down, then the Sanskrit for sits.
\end{quote}

\subsection*{You'll want to know: \sk{बिभेति}}
\textbf{\sk{बिभेति}} — \emph{bibheti} — \textquotedblleft{}is afraid\textquotedblright{}.

\begin{grammarlens}[title={What you've built}]
One more word for this chapter.
\end{grammarlens}

\subsection*{Your turn}
\begin{itemize}
\raggedright
  \item \emph{Say it:} \emph{bibheti}
  \item \emph{Say it:} \emph{bibheti}, once more
  \item \emph{Say it:} say \emph{bibheti}
  \item \emph{Recall:} say the Sanskrit for lies down, then the Sanskrit for sits, then say \emph{bibheti} again
\end{itemize}

\begin{tcolorbox}[breakable,colback=teal!4,colframe=teal!35!black,title={Before you move on}]
What does \sk{बिभेति} mean? (\textquotedblleft{}Is afraid\textquotedblright{}.) Say it once more.
\end{tcolorbox}
\section[\texorpdfstring{\sk{ध्यायति} (dhyāyati)}{ध्यायति (dhyāyati)}]{\sk{ध्यायति} (dhyāyati) — meditates, thinks of}

\label{lesson:SA-C143-dhyayati}



\begin{quote}
Before the new one: say the Sanskrit for sits, then the Sanskrit for is afraid.
\end{quote}

\subsection*{You'll want to know: \sk{ध्यायति}}
\textbf{\sk{ध्यायति}} — \emph{dhyāyati} — \textquotedblleft{}meditates, thinks of\textquotedblright{}.

\begin{grammarlens}[title={What you've built}]
One more word for this chapter.
\end{grammarlens}

\subsection*{Your turn}
\begin{itemize}
\raggedright
  \item \emph{Say it:} \emph{dhyāyati}
  \item \emph{Say it:} \emph{dhyāyati}, once more
  \item \emph{Say it:} say \emph{dhyāyati}
  \item \emph{Recall:} say the Sanskrit for sits, then the Sanskrit for is afraid, then say \emph{dhyāyati} again
\end{itemize}

\begin{tcolorbox}[breakable,colback=teal!4,colframe=teal!35!black,title={Before you move on}]
What does \sk{ध्यायति} mean? (\textquotedblleft{}Meditates, thinks of\textquotedblright{}.) Say it once more.
\end{tcolorbox}
